% =============================================================================
% Tabla de Conexiones: ESP32-S3 <-> Pantalla TFT 1.69" (ST7789V3)
% Proyecto: Prueba_TFT_S3
% Fecha: 2026-07-27
% =============================================================================
\documentclass[11pt,a4paper]{article}

% --- Paquetes ---
\usepackage[spanish,es-noshorthands]{babel}
\usepackage[utf8]{inputenc}
\usepackage[T1]{fontenc}
\usepackage{geometry}
\usepackage{booktabs}
\usepackage{array}
\usepackage{xcolor}
\usepackage{colortbl}
\usepackage{graphicx}
\usepackage{siunitx}
\usepackage{amsmath}
\usepackage{fancyhdr}
\usepackage{hyperref}

% --- Geometría de página ---
\geometry{margin=2.5cm}

% --- Colores personalizados ---
\definecolor{azulESP}{HTML}{1565C0}
\definecolor{verdeTFT}{HTML}{2E7D32}
\definecolor{grisFondo}{HTML}{F5F5F5}
\definecolor{rojoVCC}{HTML}{D32F2F}
\definecolor{negroGND}{HTML}{212121}

% --- Encabezado y pie ---
\pagestyle{fancy}
\fancyhf{}
\fancyhead[L]{\footnotesize Proyecto Prueba\_TFT\_S3}
\fancyhead[R]{\footnotesize ESP32-S3 + TFT 1.69''}
\fancyfoot[C]{\thepage}
\renewcommand{\headrulewidth}{0.4pt}

% --- Inicio del documento ---
\begin{document}

\begin{center}
    {\LARGE\bfseries\color{azulESP} Tabla de Conexiones}\\[8pt]
    {\large ESP32-S3-DevKitC-1 $\longleftrightarrow$ Pantalla TFT 1.69'' (ST7789V3)}\\[4pt]
    {\normalsize\color{gray} Resolución: 240$\times$280 px (RGB) --- Interfaz: SPI}\\[12pt]
\end{center}

% =============================================================================
% TABLA PRINCIPAL DE CONEXIONES
% =============================================================================
\section{Conexiones de Hardware}

\begin{table}[h!]
\centering
\renewcommand{\arraystretch}{1.4}
\begin{tabular}{
    >{\centering\arraybackslash}m{1.2cm}
    >{\bfseries\centering\arraybackslash}m{1.5cm}
    >{\centering\arraybackslash}m{1cm}
    >{\bfseries\centering\arraybackslash}m{2cm}
    m{5.5cm}
}
\toprule
\rowcolor{azulESP!15}
\textbf{N\textsuperscript{o}} & \textbf{Pin TFT} & \textbf{$\rightarrow$} & \textbf{Pin ESP32-S3} & \textbf{Función / Notas} \\
\midrule

1 & \textcolor{negroGND}{GND} & --- & GND & Masa o tierra común. Conectar a cualquier pin GND de la placa. \\[2pt]

\rowcolor{grisFondo}
2 & \textcolor{rojoVCC}{VCC} & --- & 5V (VBUS) & Alimentación del módulo. Conectar al pin de \SI{5}{\volt} (no a \SI{3.3}{\volt}) ya que el regulador U6 de la placa TFT requiere un voltaje de entrada superior a \SI{3.3}{\volt}. \\[2pt]

3 & SCL & --- & GPIO 12 & Señal de reloj del bus SPI (SPI Clock / SCK). \\[2pt]

\rowcolor{grisFondo}
4 & SDA & --- & GPIO 11 & Datos del bus SPI (MOSI --- Master Out Slave In). \\[2pt]

5 & RES & --- & GPIO 14 & Señal de reinicio (Reset) de la pantalla. Activo en nivel bajo. \\[2pt]

\rowcolor{grisFondo}
6 & DC & --- & GPIO 9 & Selección de Datos/Comandos (Data/Command). HIGH = dato, LOW = comando. \\[2pt]

7 & CS & --- & GPIO 10 & Selección de chip (Chip Select). Activo en nivel bajo. \\[2pt]

\rowcolor{grisFondo}
8 & BLK & --- & GPIO 21 & Control de retroiluminación (Backlight). HIGH = encendida. Alternativa: conectar directo a \SI{3.3}{\volt} para brillo constante al máximo. \\[2pt]

\bottomrule
\end{tabular}
\caption{Conexiones entre el ESP32-S3 y la pantalla TFT 1.69'' ST7789V3.}
\label{tab:conexiones}
\end{table}

% =============================================================================
% NOTAS IMPORTANTES
% =============================================================================
\section{Notas Importantes}

\begin{enumerate}
    \item \textbf{Alimentación (VCC):} El módulo TFT tiene un regulador de voltaje integrado (componente U6 en la placa). Este regulador necesita un voltaje de entrada mayor a \SI{3.3}{\volt} para funcionar correctamente. Por eso el pin VCC se conecta al pin de \SI{5}{\volt} del ESP32-S3. Los pines de datos operan a \SI{3.3}{\volt} de forma segura.
    
    \item \textbf{Puerto SPI:} El ESP32-S3 requiere la bandera \texttt{USE\_FSPI\_PORT=1} en la configuración de la librería TFT\_eSPI para utilizar el puerto SPI2 (FSPI). Sin esta configuración, el microcontrolador produce un error de tipo \textit{Guru Meditation Error (StoreProhibited)} al intentar inicializar la pantalla.
    
    \item \textbf{Velocidad SPI:} Se utiliza una frecuencia de reloj de \SI{27}{\mega\hertz} (\texttt{SPI\_FREQUENCY=27000000}), que ofrece un buen equilibrio entre velocidad de refresco y estabilidad de la señal.
    
    \item \textbf{Configuración del controlador:} Las banderas \texttt{CGRAM\_OFFSET=1}, \texttt{TFT\_INVERSION\_ON=1} y \texttt{TFT\_RGB\_ORDER=TFT\_RGB} son necesarias para que los colores y la memoria gráfica se mapeen correctamente en pantallas de 240$\times$280 píxeles.
\end{enumerate}

% =============================================================================
% CONFIGURACIÓN DE SOFTWARE
% =============================================================================
\section{Configuración de Software (PlatformIO)}

La siguiente configuración se define en el archivo \texttt{platformio.ini}:

\begin{table}[h!]
\centering
\renewcommand{\arraystretch}{1.3}
\begin{tabular}{ll}
\toprule
\rowcolor{azulESP!15}
\textbf{Parámetro} & \textbf{Valor} \\
\midrule
Plataforma              & \texttt{espressif32} \\
\rowcolor{grisFondo}
Placa                   & \texttt{esp32-s3-devkitc-1} \\
Framework               & \texttt{Arduino} \\
\rowcolor{grisFondo}
Librería                & \texttt{bodmer/TFT\_eSPI} \\
Controlador             & \texttt{ST7789\_DRIVER} \\
\rowcolor{grisFondo}
Resolución              & 240 $\times$ 280 px \\
Puerto SPI              & FSPI (SPI2) \\
\rowcolor{grisFondo}
Frecuencia SPI          & \SI{27}{\mega\hertz} \\
Velocidad serial        & 115200 baud \\
\bottomrule
\end{tabular}
\caption{Parámetros de configuración del proyecto en PlatformIO.}
\label{tab:config}
\end{table}

\end{document}
